\documentclass[11pt,a4paper]{article}
\usepackage[utf8]{inputenc}
\usepackage[margin=1in]{geometry}
\usepackage{graphicx}
\usepackage{amsmath}
\usepackage{amssymb}
\usepackage{hyperref}
\usepackage{xcolor}
\usepackage{fancyhdr}
\usepackage{listings}
\usepackage{booktabs}
\usepackage{float}

% Header and Footer
\pagestyle{fancy}
\fancyhf{}
\rhead{Aircraft Predictive Maintenance}
\lhead{Technical Report}
\cfoot{\thepage}

% Title
\title{\textbf{Aircraft Predictive Maintenance System}\\
\large A Machine Learning Approach to Component Failure Prediction}
\author{Civil Aviation Project}
\date{\today}

\begin{document}

\maketitle
\tableofcontents
\newpage

% ============================================
\section{Executive Summary}
% ============================================

This report documents the development of a predictive maintenance system for aircraft components using machine learning. The system predicts component failures based on 12 sensor inputs with 98.76\% ROC-AUC accuracy, enabling proactive maintenance scheduling. The model achieves 100\% recall on critical failures, ensuring no maintenance needs are missed while minimizing false alarms at 62.9\% precision.

\textbf{Key Metrics:}
\begin{itemize}
    \item ROC-AUC Score: 98.76\%
    \item Recall: 100\% (all failures detected)
    \item Precision: 62.9\% (acceptable false positive rate)
    \item Model Type: Logistic Regression with SMOTE
    \item Input Features: 12 sensor measurements
\end{itemize}

% ============================================
\section{Problem Understanding}
% ============================================

\subsection{Business Context}
Aircraft component failures lead to:
\begin{itemize}
    \item Unplanned maintenance costs (30-40\% higher than scheduled)
    \item Flight delays and cancellations
    \item Safety risks if failures go undetected
    \item Reduced aircraft availability
\end{itemize}

\subsection{Problem Statement}
\textbf{Objective:} Predict which aircraft components will fail in the next maintenance cycle, allowing maintenance teams to replace or service components proactively.

\textbf{Classification Task:} Binary classification
\begin{itemize}
    \item \textbf{Class 0 (Negative):} Component is operating normally
    \item \textbf{Class 1 (Positive):} Component will fail soon and requires maintenance
\end{itemize}

\subsection{Success Criteria}
\begin{enumerate}
    \item \textbf{High Recall:} Minimize missed failures (safety priority)
    \item \textbf{Reasonable Precision:} Avoid excessive false alarms (cost priority)
    \item \textbf{Interpretability:} Maintenance teams must understand predictions
    \item \textbf{Scalability:} Handle multiple aircraft and component types
\end{enumerate}

% ============================================
\section{Data Exploration and Insights}
% ============================================

\subsection{Dataset Overview}
\begin{table}[H]
\centering
\begin{tabular}{|l|r|}
\toprule
\textbf{Characteristic} & \textbf{Value} \\
\midrule
Total Records & \textgreater 1,000 \\
Features & 12 sensor measurements \\
Target Variable & Binary (0 = Normal, 1 = Failure) \\
Class Distribution & Imbalanced (majority normal) \\
Missing Values & None after preprocessing \\
\bottomrule
\end{tabular}
\caption{Dataset Characteristics}
\end{table}

\subsection{Sensor Features}
The dataset includes 12 sensor measurements from aircraft monitoring systems:

\begin{itemize}
    \item \textbf{Temperature Sensors:} Engine, hydraulic, and bearing temperatures
    \item \textbf{Vibration Sensors:} Component vibration measurements at different frequencies
    \item \textbf{Pressure Sensors:} Oil pressure, fuel pressure, cabin pressure
    \item \textbf{Efficiency Metrics:} Fuel consumption, operational hours
\end{itemize}

\subsection{Class Imbalance Problem}
\textbf{Challenge:} Only 10-15\% of records represent actual failures (positive class), while 85-90\% represent normal operations (negative class).

\textbf{Impact:} Standard models tend to:
\begin{itemize}
    \item Predict ``normal'' for most cases (high accuracy, low recall on failures)
    \item Miss critical failures due to class imbalance
    \item Require specialized techniques to handle minority class
\end{itemize}

\subsection{Key Data Insights}
\begin{enumerate}
    \item \textbf{Separated Distributions:} Failed components show distinct sensor patterns
    \item \textbf{Temporal Trends:} Sensor values gradually increase before failure
    \item \textbf{Multi-sensor Correlation:} Failures often involve multiple sensors simultaneously
    \item \textbf{Outliers:} Some extreme values indicate imminent failure (useful signals)
\end{enumerate}

% ============================================
\section{Feature Engineering and Preprocessing}
% ============================================

\subsection{Data Cleaning}
\begin{itemize}
    \item Removed duplicate records
    \item Handled outliers using IQR method (kept extreme values as failure indicators)
    \item Verified no missing values
    \item Removed non-predictive features
\end{itemize}

\subsection{Feature Scaling}
Applied \textbf{StandardScaler} to normalize all 12 sensor features:
\[
x_{\text{scaled}} = \frac{x - \mu}{\sigma}
\]

\textbf{Rationale:}
\begin{itemize}
    \item Logistic Regression performs better with scaled features
    \item Prevents high-magnitude sensors from dominating predictions
    \item Ensures consistent feature importance interpretation
\end{itemize}

\subsection{Handling Class Imbalance}
Applied \textbf{SMOTE (Synthetic Minority Over-sampling Technique)} to training set:
\begin{itemize}
    \item Creates synthetic samples of minority class (failures)
    \item Balances training data to 50\%-50\% distribution
    \item Prevents model from ignoring critical failure patterns
    \item Applied \textbf{only} to training set (not validation/test)
\end{itemize}

\textbf{Key Decision:} SMOTE only on training data ensures unbiased evaluation metrics.

\subsection{Train-Test Split}
\begin{itemize}
    \item \textbf{Training Set:} 70\% of data (SMOTE applied)
    \item \textbf{Test Set:} 30\% of data (original distribution)
    \item Stratified split to maintain class proportions
\end{itemize}

% ============================================
\section{Model Selection and Evaluation}
% ============================================

\subsection{Model Choice: Logistic Regression}

\textbf{Why Logistic Regression?}
\begin{table}[H]
\centering
\begin{tabular}{|l|c|c|}
\toprule
\textbf{Criterion} & \textbf{Logistic Regression} & \textbf{Notes} \\
\midrule
Interpretability & Excellent & Feature coefficients are clear \\
Training Speed & Very Fast & Suitable for real-time deployment \\
Scalability & Excellent & Handles large datasets easily \\
Probability Outputs & Native & Confidence scores for maintenance \\
Explainability & High & Maintenance teams understand predictions \\
\bottomrule
\end{tabular}
\caption{Logistic Regression Advantages}
\end{table}

Compared alternatives (Random Forest, SVM) but selected Logistic Regression for:
\begin{enumerate}
    \item \textbf{Production Simplicity:} Easy to deploy and maintain
    \item \textbf{Explainability:} Maintenance engineers need to trust predictions
    \item \textbf{Real-time Performance:} Fast inference for continuous monitoring
    \item \textbf{Calibration:} Probability outputs align with maintenance planning
\end{enumerate}

\subsection{Model Performance Metrics}

\subsubsection{Classification Metrics}
\begin{table}[H]
\centering
\begin{tabular}{|l|r|p{5cm}|}
\toprule
\textbf{Metric} & \textbf{Value} & \textbf{Interpretation} \\
\midrule
Accuracy & 96.2\% & Overall correctness \\
Precision & 62.9\% & Of predicted failures, 62.9\% are real \\
Recall (Sensitivity) & 100\% & \textbf{All failures detected} \\
Specificity & 92.1\% & Normal components correctly identified \\
F1-Score & 0.773 & Balanced precision-recall measure \\
ROC-AUC & 98.76\% & Excellent ranking of predictions \\
\bottomrule
\end{tabular}
\caption{Model Performance on Test Set}
\end{table}

\subsubsection{Interpretation}
\begin{itemize}
    \item \textbf{Perfect Recall (100\%):} Zero false negatives---no failures missed. Safety ensured.
    \item \textbf{Good Precision (62.9\%):} 37.1\% false positives acceptable for predictive maintenance (preventive cost \textless preventive benefits)
    \item \textbf{ROC-AUC 98.76\%:} Excellent discrimination between normal and failing components
\end{itemize}

\subsection{Decision Threshold}
Current threshold: 0.5 probability
\begin{itemize}
    \item \textbf{Consequence If Changed:}
    \begin{itemize}
        \item Lowering threshold → More false positives (higher maintenance costs)
        \item Raising threshold → More false negatives (increased failure risk)
    \end{itemize}
    \item \textbf{Recommendation:} Keep at 0.5 (balanced approach for maintenance planning)
\end{itemize}

% ============================================
\section{Key Challenges and Assumptions}
% ============================================

\subsection{Challenges Encountered}

\subsubsection{1. Class Imbalance}
\textbf{Challenge:} Only 10-15\% of data represents failures.
\begin{itemize}
    \item Direct models would predict ``all normal'' and achieve 90\% accuracy (useless)
    \item Solution: SMOTE balancing during training
\end{itemize}

\subsubsection{2. Feature Correlation}
\textbf{Challenge:} Some sensor measurements are highly correlated.
\begin{itemize}
    \item Could inflate feature importance
    \item Solution: Acceptable for Logistic Regression (linear relationships preserved)
\end{itemize}

\subsubsection{3. Non-stationary Patterns}
\textbf{Challenge:} Aircraft components degrade differently over time.
\begin{itemize}
    \item Historical patterns may not apply to new aircraft types
    \item Solution: Regular model retraining with new maintenance data
\end{itemize}

\subsubsection{4. Domain Knowledge Integration}
\textbf{Challenge:} Model trained on data; lacks expert maintenance rules.
\begin{itemize}
    \item Solution: Maintenance team reviews questionable predictions
    \item Model works as decision support, not autonomous system
\end{itemize}

\subsection{Key Assumptions}

\begin{enumerate}
    \item \textbf{Assumption 1:} Historic failure patterns will repeat
    \begin{itemize}
        \item Validity: High (aircraft systems are standardized)
        \item Risk: Mitigated by quarterly retraining
    \end{itemize}
    
    \item \textbf{Assumption 2:} Sensor data quality remains consistent
    \begin{itemize}
        \item Validity: Medium (sensors can degrade)
        \item Risk: Mitigation: Input validation and anomaly detection
    \end{itemize}
    
    \item \textbf{Assumption 3:} Features capture sufficient failure information
    \begin{itemize}
        \item Validity: High (12 comprehensive sensors)
        \item Risk: Acceptable (domain experts validated feature set)
    \end{itemize}
    
    \item \textbf{Assumption 4:} Test data distribution matches production
    \begin{itemize}
        \item Validity: Medium (new aircraft types may differ)
        \item Risk: Mitigated by continuous monitoring of model drift
    \end{itemize}
\end{enumerate}

% ============================================
\section{Deployment and System Architecture}
% ============================================

\subsection{Technology Stack}
\begin{itemize}
    \item \textbf{Backend:} FastAPI (Python) with REST API endpoints
    \item \textbf{Frontend:} React.js for user interface
    \item \textbf{Deployment:} Railway.app (cloud hosting)
    \item \textbf{Model Serialization:} joblib (best\_model.pkl)
\end{itemize}

\subsection{API Endpoints}
\begin{itemize}
    \item \textbf{POST /predict} - Single component prediction
    \item \textbf{POST /predict\_batch} - Batch predictions for multiple components
    \item \textbf{GET /model\_info} - Model metrics and metadata
    \item \textbf{GET /health} - System health check
\end{itemize}

\subsection{Live URL}
\textbf{Production API:} \url{https://civil-aviationprojectaircraftmaintenance-production.up.railway.app}

% ============================================
\section{Model Interpretability: SHAP Explanations}
% ============================================

\subsection{Why Explainability Matters}

Maintenance teams must understand \textbf{WHY} the model predicts failure, not just the prediction itself. A black-box model cannot be trusted for safety-critical decisions.

**Problem:** Logistic Regression provides predictions but no clear explanations
- Maintenance question: ``Why is this component flagged as at-risk?''
- Without explanation: Cannot verify if prediction makes sense
- Impact: Team cannot trust or act on automated predictions

\subsection{SHAP (SHapley Additive exPlanations)}

SHAP uses game theory to calculate each feature's contribution to predictions:

\[
\text{Model Output} = \text{Base Value} + \sum_{i=1}^{12} \text{Feature}_i \text{ Contribution}
\]

\textbf{Example Interpretation:}
\begin{verbatim}
Base Failure Probability:          15%
+ High Temperature (sensor_1):     +25%  (RED: pushes toward failure)
- Normal Vibration (sensor_3):      -5%  (BLUE: pushes toward normal)
+ High Pressure (sensor_5):        +18%  (RED: pushes toward failure)
────────────────────────────────────────
= Final Prediction:                53%   (MAINTENANCE REQUIRED)
\end{verbatim}

\subsection{SHAP Value Interpretation}

\textbf{Red Arrows (Risk Factors):} Increase failure probability
- Example: High engine temperature
- Action: Flag for maintenance

\textbf{Blue Arrows (Protective Factors):} Decrease failure probability
- Example: Normal vibration levels
- Action: Reassure that component is stable

\subsection{Example: Explaining High-Risk Prediction}

\textbf{Component Status Report:}
\begin{itemize}
    \item \textbf{Prediction:} Failure Risk (1)
    \item \textbf{Probability:} 76\% (+0.76)
    \item \textbf{Confidence:} High
\end{itemize}

\textbf{Top Risk Factors (driving toward failure):}
\begin{table}[H]
\centering
\begin{tabular}{|l|r|p{4cm}|}
\toprule
\textbf{Feature} & \textbf{SHAP Value} & \textbf{Interpretation} \\
\midrule
sensor\_1 (Temp) & +0.34 & HIGH: Extreme temperature \\
sensor\_5 (Pressure) & +0.28 & HIGH: Elevated pressure \\
sensor\_7 (Vibration) & +0.22 & MEDIUM: Excessive vibration \\
sensor\_11 (Hours) & +0.12 & LOW: High operational hours \\
\bottomrule
\end{tabular}
\end{table}

\textbf{Recommendation for Maintenance Team:}
\begin{itemize}
    \item Focus on cooling system (high temperature)
    \item Check pressure relief valve
    \item Inspect bearing alignment (vibration)
    \item Schedule maintenance within next 100 flight hours
\end{itemize}

% ============================================
\section{Data Drift Detection and Monitoring}
% ============================================

\subsection{The Data Drift Problem}

Over time, the distribution of production data can shift from training data, causing model performance to degrade:

\textbf{Example Timeline:}
\begin{itemize}
    \item Month 1 (Deployment): Model accuracy = 96\% ✓
    \item Month 3: Model accuracy = 91\% ⚠️
    \item Month 6: Model accuracy = 82\% ❌ (Unacceptable)
\end{itemize}

\textbf{Root Cause:} Aircraft sensor distributions changed, but model wasn't retrained. Model saw old patterns but new data is different.

\subsection{Types of Data Drift}

\begin{table}[H]
\centering
\begin{tabular}{|l|p{3cm}|p{3cm}|p{3cm}|}
\toprule
\textbf{Type} & \textbf{Cause} & \textbf{Impact} & \textbf{Detection} \\
\midrule
Covariate Drift & Input feature distribution changes & Model performance degrades & Statistical tests on features \\
Label Drift & Target variable distribution shifts & Predictions become biased & Compare model output vs actual \\
Concept Drift & Feature-target relationship changes & Model assumptions violated & Track accuracy over time \\
\bottomrule
\end{tabular}
\end{table}

\subsection{Example: Temperature Sensor Drift}

\textbf{Training Data (2025):}
\begin{itemize}
    \item Average Engine Temperature: 45°C
    \item Standard Deviation: 5°C
    \item Range: 35-55°C
\end{itemize}

\textbf{Production Data (6 months later):}
\begin{itemize}
    \item Average Engine Temperature: 62°C
    \item Standard Deviation: 8°C
    \item Range: 50-80°C
\end{itemize}

\textbf{Problem:} Model trained on 35-55°C range, never saw temperatures $\geq$ 50°C. New production data consistently above 50°C → Predictions unreliable!

\subsection{Drift Detection Methods}

\subsubsection{Kolmogorov-Smirnov Test}

Statistical hypothesis test comparing two distributions:

\[
D = \max |F_{\text{train}}(x) - F_{\text{prod}}(x)|
\]

\textbf{Decision Rule:}
\begin{itemize}
    \item p-value $> 0.05$: NO drift (distributions similar) ✓
    \item p-value $\leq 0.05$: DRIFT DETECTED ⚠️
    \item p-value $< 0.01$: CRITICAL DRIFT ❌
\end{itemize}

\subsubsection{Population Stability Index (PSI)}

Measures shift in feature distribution:

\[
\text{PSI} = \sum_{i} \left( \% \text{prod}_i - \% \text{train}_i \right) \times \ln\left(\frac{\% \text{prod}_i}{\% \text{train}_i}\right)
\]

\textbf{Interpretation:}
\begin{itemize}
    \item PSI $< 0.1$: No significant drift
    \item PSI $0.1-0.25$: Moderate drift (monitor)
    \item PSI $> 0.25$: Severe drift (retrain immediately)
\end{itemize}

\subsection{Monitoring and Alerting}

\textbf{Recommended Monitoring Schedule:}

\begin{table}[H]
\centering
\begin{tabular}{|l|l|l|}
\toprule
\textbf{Frequency} & \textbf{Check} & \textbf{Action} \\
\midrule
Daily & Model predictions (count) & Ensure system is active \\
Weekly & Drift tests on all features & Log results, set alerts \\
Monthly & Model accuracy on new data & Compare vs baseline \\
Quarterly & Full model retraining & Update with recent data \\
\bottomrule
\end{tabular}
\end{table}

\subsection{Retraining Pipeline}

When drift is detected, follow this workflow:

\textbf{Step 1: Alert Maintenance Team}
\begin{verbatim}
"Model accuracy degraded from 96% to 82%.
Using backup maintenance rules until retraining complete."
\end{verbatim}

\textbf{Step 2: Collect New Data}
\begin{itemize}
    \item Gather recent production predictions
    \item Get maintenance outcomes (actual failures)
    \item Combine with original training data
\end{itemize}

\textbf{Step 3: Retrain Model}
\begin{itemize}
    \item Balance classes with SMOTE
    \item Apply StandardScaler
    \item Train new Logistic Regression model
    \item Validate on recent held-out test set
\end{itemize}

\textbf{Step 4: Validate}
\begin{itemize}
    \item Recall $\geq 95\%$ (safety requirement)
    \item Precision $\geq 60\%$ (cost constraint)
    \item Compare against old model on same data
\end{itemize}

\textbf{Step 5: Deploy}
\begin{itemize}
    \item Replace old model with new version
    \item Update model artifact file
    \item Commit to GitHub (triggers auto-deployment)
    \item Monitor closely for next 2 weeks
\end{itemize}

% ============================================
\section{Recommendations and Future Work}
% ============================================

\subsection{Short-term (Next 3 months)}
\begin{enumerate}
    \item Deploy interactive dashboard for maintenance team
    \item Establish feedback loop for prediction corrections
    \item Implement automated alerts for high-risk components
\end{enumerate}

\subsection{Medium-term (3-6 months)}
\begin{enumerate}
    \item Collect additional failure data for model refinement
    \item Implement model performance monitoring
    \item Extend to other component types
\end{enumerate}

\subsection{Long-term (6+ months)}
\begin{enumerate}
    \item Explore deep learning for temporal patterns
    \item Integrate with maintenance scheduling system
    \item Expand to predictive cost estimation
\end{enumerate}

% ============================================
\section{Conclusion}
% ============================================

The developed predictive maintenance system achieves excellent performance with 100\% recall and 98.76\% ROC-AUC, ensuring no critical failures are missed. The use of Logistic Regression provides interpretability and fast inference, making it suitable for real-time deployment in aircraft maintenance operations.

The system is now live and accessible via REST API. Maintenance teams can use predictions to schedule proactive maintenance, reducing costs and improving safety. With continuous monitoring and regular retraining, this system will provide long-term value to aircraft maintenance operations.

\end{document}
